\section{Introduction}
\label{sec:intro}

\paragraph{Abstract.}
Vision-Language Models (\vlm{}s) routinely accept user-supplied images alongside benign text prompts. We ask whether a single image, perceptually indistinguishable from a real photo, can carry hidden instructions that the \vlm{} obeys at inference time. We build \textsc{VisInject}, a three-stage pipeline that (i) optimises a universal adversarial image against a small ensemble of open white-box \vlm{}s, (ii) transports the resulting attack signal onto arbitrary clean photos via a pretrained \textsc{AnyAttack} decoder \citep{zhang2025anyattack} under an \Linf{} budget of $\eps = 16/255$, and (iii) evaluates 6{,}615 (clean, adversarial) response pairs along two independent axes: \emph{Output Affected} (was the answer disturbed?) and \emph{Target Injected} (did the chosen payload appear?). Across $147$ experimental runs covering seven target phrases, three white-box configurations, and seven test images, we measure a disruption rate of $\approx 66\%$ but a literal-injection rate of only $0.227\%$ ($15/6{,}615$). Two confirmed injections, three partial cases, and five weak cases concentrate where the image semantics already invite text transcription. A manual transfer test on the strongest case fails on GPT-4o, which recognises the noise as ``distortion'' and recovers the underlying content.

\paragraph{Motivation.}
Image inputs are now an ordinary part of how users interact with hosted assistants, autonomous agents, and tool-using language models. Prior work on indirect prompt injection \citep{greshake2023indirect} and on adversarial robustness of multimodal models \citep{schlarmann2023robustness, bagdasaryan2023imagessounds} suggests that the visual channel is exploitable in principle. The practical question for a course-scale study is which parts of that attack surface actually deliver a payload, under realistic perceptual and threat-model constraints.

\paragraph{Contributions.}
This report makes three concrete contributions:
\begin{itemize}
  \item \textbf{A reproducible attack pipeline.} \textsc{VisInject} composes a multi-prompt PGD optimisation \citep{madry2018pgd} with an off-the-shelf \textsc{AnyAttack} decoder so that a \emph{single} universal image can be turned into many adversarial photos at $\Linf = 16/255$ and PSNR $\approx 25.2$~dB.
  \item \textbf{A dual-dimension evaluation that separates disruption from injection.} An earlier judge-style evaluation conflated ``did the answer change?'' with ``did the payload appear?''~--~we show this inflates reported success from $0.4\%$ to $50.5\%$. Our second-version evaluation reports both numbers separately.
  \item \textbf{An empirical map of where injection actually lands.} The two confirmed injections both occur on a code screenshot with a URL payload; the partial and weak cases follow the same pattern of semantic compatibility between image and target. Architecture matters far more than scale --- BLIP-2 is fully immune across all $2{,}205$ pairs while every Qwen / DeepSeek model under test is disrupted on $\geq 98\%$ of pairs.
\end{itemize}

The remainder of the report develops the threat model (\S\ref{sec:threat}), the three-stage pipeline (\S\ref{sec:method}), the experimental setup (\S\ref{sec:experiments}), aggregate results (\S\ref{sec:results}), case studies (\S\ref{sec:cases}), and a discussion of why the gap between disruption and injection is what it is (\S\ref{sec:discussion}).
